\documentclass[10pt,aspectratio=169]{beamer}
% Preámbulo modular para presentaciones Beamer (Modelación Estadística)
% Diseño y paleta institucional del Tecnológico de Monterrey

\documentclass[aspectratio=169,xcolor={dvipsnames,table}]{beamer}

\usepackage[utf8]{inputenc}
\usepackage[spanish,mexico]{babel}
\usepackage{amsmath,amssymb,amsfonts,mathtools}
\usepackage{graphicx}
\usepackage{listings}
\usepackage{booktabs}
\usepackage{tikz}
\usetikzlibrary{matrix,arrows,decorations.pathmorphing,shapes.geometric,positioning}

% Paleta Institucional Tecnológico de Monterrey
\definecolor{TecAzul}{HTML}{0039A6}       % Azul TEC: color primario institucional
\definecolor{TecAzulOscuro}{HTML}{1A2E51} % Azul marino oscuro
\definecolor{TecRojo}{HTML}{EC2661}       % Rojo de acento (paleta miTec)
\definecolor{TecGris}{HTML}{646464}       % Gris neutro para comentarios y subtítulos
\definecolor{TecFondoBloque}{HTML}{F4F6F9}% Fondo suave para bloques

% Configuración de Tema Beamer
\usetheme{Boadilla}
\usecolortheme[named=TecAzulOscuro]{structure}
\setbeamercolor{alerted text}{fg=TecRojo}
\setbeamercolor{palette primary}{bg=TecAzulOscuro,fg=white}
\setbeamercolor{palette secondary}{bg=TecAzul,fg=white}
\setbeamercolor{palette tertiary}{bg=TecAzulOscuro!80!black,fg=white}
\setbeamercolor{title}{fg=TecAzulOscuro,bg=TecFondoBloque}
\setbeamercolor{frametitle}{fg=TecAzulOscuro,bg=TecFondoBloque}

% Bloques personalizados elegantes
\setbeamercolor{block title}{bg=TecAzulOscuro,fg=white}
\setbeamercolor{block body}{bg=TecFondoBloque,fg=black}
\setbeamercolor{block title alerted}{bg=TecRojo,fg=white}
\setbeamercolor{block body alerted}{bg=TecRojo!8,fg=black}
\setbeamercolor{block title example}{bg=TecAzul,fg=white}
\setbeamercolor{block body example}{bg=TecAzul!8,fg=black}

% Redondear bordes y añadir sombra sutil
\setbeamertemplate{blocks}[rounded][shadow=true]
\setbeamertemplate{navigation symbols}{}

% Pie de página limpio con número de diapositiva
\setbeamertemplate{footline}{
  \leavevmode%
  \hbox{%
  \begin{beamercolorbox}[wd=.7\paperwidth,ht=2.25ex,dp=1ex,left]{author in head/foot}%
    \hspace*{2ex}\insertshortauthor~~(\insertshortinstitute) \hfill \insertshorttitle
  \end{beamercolorbox}%
  \begin{beamercolorbox}[wd=.3\paperwidth,ht=2.25ex,dp=1ex,right]{date in head/foot}%
    \insertshortdate{}\hspace*{2em}
    \insertframenumber{} / \inserttotalframenumber\hspace*{2ex}
  \end{beamercolorbox}}%
  \vskip0pt%
}

% Configuración de Listings de Python para visualización pedagógica de código
\lstset{
    language=Python,
    basicstyle=\ttfamily\footnotesize,
    keywordstyle=\color{TecAzulOscuro}\bfseries,
    stringstyle=\color{TecRojo},
    commentstyle=\color{TecGris}\itshape,
    showstringspaces=false,
    breaklines=true,
    frame=lines,
    rulecolor=\color{TecAzulOscuro!30},
    backgroundcolor=\color{TecFondoBloque},
    aboveskip=0.5em,
    belowskip=0.5em,
    literate={á}{{\'a}}1 {é}{{\'e}}1 {í}{{\'i}}1 {ó}{{\'o}}1 {ú}{{\'u}}1
             {Á}{{\'A}}1 {É}{{\'E}}1 {Í}{{\'I}}1 {Ó}{{\'O}}1 {Ú}{{\'U}}1
             {ñ}{{\~n}}1 {Ñ}{{\~N}}1 {¿}{{?`}}1 {¡}{{!`}}1
}

% Metadatos institucionales por defecto
\title[Modelación Estadística]{Modelación Estadística para Ciencia de Datos e Ingeniería}
\author[J. Castillo Colmenares]{Juliho David Castillo Colmenares}
\institute[Tec de Monterrey]{Tecnológico de Monterrey \\ Escuela de Ingeniería y Ciencias}
\date{\today}

% Comandos y macros estadísticas compartidas para las presentaciones Beamer (Metropolis)

% Conjuntos numéricos básicos
\newcommand{\R}{\mathbb{R}}
\newcommand{\N}{\mathbb{N}}
\newcommand{\Q}{\mathbb{Q}}
\newcommand{\Z}{\mathbb{Z}}
\newcommand{\set}[1]{\left\{ #1 \right\}}
\newcommand{\abs}[1]{\left|#1\right|}
\newcommand{\norm}[1]{\left\| #1 \right\|}

% Comandos estadísticos y probabilísticos del curso
\newcommand{\Var}{\operatorname{Var}}
\newcommand{\cov}{\operatorname{Cov}}
\newcommand{\corr}{\rho}
\newcommand{\comb}[2]{\begin{pmatrix} #1 \\ #2 \end{pmatrix}}
\newcommand{\s}{\sigma}
\newcommand{\particion}{\mathfrak{P}}
\newcommand{\card}[1]{\#\left( #1 \right)}
\newcommand{\seq}[3]{\set{#1}_{#2}^{#3}}
\newcommand{\E}{\mathbb{E}}
\newcommand{\Prob}{\mathbb{P}}

% Entornos pedagógicos y cajas en español académico natural
\newenvironment{discusion}{%
  \begin{alertblock}{Pregunta de Discusión}%
}{%
  \end{alertblock}%
}

\newenvironment{reflexion}{%
  \begin{alertblock}{Reflexión para el Aula}%
}{%
  \end{alertblock}%
}

\newenvironment{paradoja}{%
  \begin{alertblock}{Paradoja de Exploración}%
}{%
  \end{alertblock}%
}

\newenvironment{motivacion}{%
  \begin{exampleblock}{Motivación: Ciencia de Datos en la Práctica}%
}{%
  \end{exampleblock}%
}

\newenvironment{retoclase}{%
  \begin{block}{Reto Rápido en Equipo}%
}{%
  \end{block}%
}


\title[Introducción a la estadística inferencial]{Sección 5.0: Introducción a la estadística inferencial}
\subtitle{De la muestra a las conclusiones sobre una población}
\author[J. Castillo Colmenares]{Juliho Castillo Colmenares}
\institute{Tecnológico de Monterrey}
\date{}

\begin{document}
\begin{frame}[plain]\vspace{-0.8cm}\titlepage\end{frame}

\begin{frame}{Hoja de ruta --- capítulo 5}
  \footnotesize
  \begin{itemize}\itemsep=0.08cm
    \item \textbf{05.00} Introducción a la estadística inferencial \emph{(hoy)}
    \item \textbf{05.01--05.02} Transformaciones y distribuciones muestrales
    \item \textbf{05.03--05.05} Chi-cuadrada, $t$ y $F$
    \item \textbf{05.08--05.09} Conceptos fundamentales y estadísticos Z/t
  \end{itemize}
  \begin{block}{Objetivo didáctico}
    Distinguir población, muestra, parámetro y estadístico; describir el proceso
    de inferencia y reconocer el papel del error muestral y del TLC.
  \end{block}
\end{frame}

\begin{frame}{¿Por qué inferir?}
  \footnotesize
  La población completa rara vez puede observarse. La nota identifica cuatro
  restricciones prácticas:
  \begin{itemize}\itemsep=0.1cm
    \item costo de medir a todos;
    \item tiempo de recolección;
    \item accesibilidad limitada;
    \item mediciones destructivas, como probar la vida útil de una bombilla.
  \end{itemize}
  \pause
  \begin{alertblock}{Idea central}
    Una muestra representativa permite generalizar, pero siempre introduce
    incertidumbre que debe cuantificarse.
  \end{alertblock}
\end{frame}

\begin{frame}{Población y muestra}
  \footnotesize
  \begin{block}{Definiciones de la sección}
    \textbf{Población:} conjunto completo de individuos, objetos o eventos con
    una característica común de interés.

    \medskip
    \textbf{Muestra:} subconjunto seleccionado para realizar el estudio.
  \end{block}
  \pause
  \begin{exampleblock}{Lectura operacional}
    La población define el alcance de la conclusión; la muestra aporta los
    datos concretos con los que se calcula.
  \end{exampleblock}
\end{frame}

\begin{frame}{Parámetro y estadístico}
  \footnotesize
  \begin{columns}[T]
    \column{0.48\textwidth}
      \begin{block}{Población}
        Parámetro fijo, generalmente desconocido: $\mu$, $\sigma$.
      \end{block}
    \column{0.48\textwidth}
      \begin{block}{Muestra}
        Estadístico calculable: $\bar{x}$, $s$.
      \end{block}
  \end{columns}
  \vspace{0.15cm}
  \begin{alertblock}{Distinción}
    El parámetro describe la población; el estadístico es una variable aleatoria
    que cambia de muestra en muestra y se usa para estimarlo.
  \end{alertblock}
\end{frame}

\begin{frame}{Dos formas de inferencia}
  \footnotesize
  \begin{enumerate}\itemsep=0.18cm
    \item \textbf{Estimación:} usar estadísticos para aproximar un parámetro.
      Puede ser puntual o mediante un intervalo de confianza.
    \item \textbf{Prueba de hipótesis:} formular una afirmación sobre un
      parámetro y decidir si la evidencia muestral la respalda o contradice.
  \end{enumerate}
  \pause
  \begin{exampleblock}{Pregunta guía}
    ¿Queremos producir un valor o rango plausible, o tomar una decisión sobre
    una afirmación específica?
  \end{exampleblock}
\end{frame}

\begin{frame}{Proceso general de inferencia}
  \footnotesize
  \begin{enumerate}\itemsep=0.12cm
    \item Definir la población y la característica de interés.
    \item Seleccionar una muestra representativa.
    \item Calcular estadísticos muestrales.
    \item Aplicar estimación o pruebas con distribuciones apropiadas.
    \item Interpretar el resultado en el contexto original.
  \end{enumerate}
  \begin{block}{Secuencia lógica}
    \centering población $\longrightarrow$ muestra $\longrightarrow$
    estadístico $\longrightarrow$ inferencia $\longrightarrow$ decisión.
  \end{block}
\end{frame}

\begin{frame}{Fuentes de error}
  \footnotesize
  \begin{columns}[T]
    \column{0.48\textwidth}
      \begin{block}{Error de muestreo}
        Diferencia aleatoria entre estadístico y parámetro. Disminuye al
        aumentar el tamaño de muestra.
      \end{block}
    \column{0.48\textwidth}
      \begin{block}{Error no muestral}
        Sesgo de selección, medición o no respuesta. No desaparece solo por
        aumentar $n$.
      \end{block}
  \end{columns}
  \begin{alertblock}{Advertencia}
    La inferencia cuantifica principalmente el error de muestreo; el diseño del
    estudio debe controlar los errores sistemáticos.
  \end{alertblock}
\end{frame}

\begin{frame}{Distribución muestral}
  \footnotesize
  Si se toman muchas muestras de tamaño $n$ y se calcula un estadístico en cada
  una, los resultados forman su \emph{distribución muestral}.
  \begin{block}{Ejemplo conceptual de la sección}
    Para cada muestra se obtiene una media $\bar X$. La distribución de todas
    esas medias describe la variabilidad de la estimación, no la variabilidad
    de las observaciones individuales.
  \end{block}
  \pause
  \begin{alertblock}{Puente}
    Esta distribución será la base de intervalos de confianza y pruebas de
    hipótesis en las secciones siguientes.
  \end{alertblock}
\end{frame}

\begin{frame}{Puente computacional 1/4: mapa de cantidades}
  \footnotesize
  \begin{tabular}{lll}
    \hline
    Objeto & Símbolo & Papel en la inferencia \\
    \hline
    Población & $\mu,\sigma$ & valores objetivo \\
    Muestra & $X_1,\ldots,X_n$ & datos observados \\
    Estadístico & $\bar X,s$ & calculables \\
    Inferencia & IC o prueba & conclusión con incertidumbre \\
    \hline
  \end{tabular}
  \vspace{0.15cm}
  La tabla resume las definiciones de la nota sin introducir un procedimiento
  externo.
\end{frame}

\begin{frame}{Puente computacional 2/4: media muestral}
  \footnotesize
  Para una muestra aleatoria con media poblacional $\mu$ y varianza $\sigma^2$,
  la cantidad que se repite entre muestras es
  \[
    \bar X=\frac{1}{n}\sum_{i=1}^{n}X_i.
  \]
  \begin{block}{Lectura de la sección}
    La media muestral es el estadístico elemental que conecta los datos con el
    parámetro $\mu$.
  \end{block}
\end{frame}

\begin{frame}{Puente computacional 3/4: variabilidad de la media}
  \footnotesize
  \[
    E(\bar X)=\mu,\qquad \operatorname{Var}(\bar X)=\frac{\sigma^2}{n}.
  \]
  \begin{itemize}\itemsep=0.13cm
    \item La media es insesgada: su centro coincide con $\mu$.
    \item Su error estándar es $\sigma/\sqrt n$.
    \item Cuadruplicar $n$ reduce a la mitad la desviación estándar de la media.
  \end{itemize}
\end{frame}

\begin{frame}{Puente computacional 4/4: TLC}
  \footnotesize
  \[
    \frac{\bar X-\mu}{\sigma/\sqrt n}\xrightarrow{d}N(0,1)
    \qquad (n\text{ grande}).
  \]
  \begin{block}{Consecuencia}
    La forma de la población original deja de dominar la distribución de la
    media cuando el tamaño de muestra es suficientemente grande.
  \end{block}
\end{frame}

\begin{frame}{Caso de lectura 1/4 --- población y muestra}
  \footnotesize
  Una medición destructiva impide estudiar todos los objetos de una población.
  La nota propone seleccionar una muestra representativa.
  \begin{alertblock}{Decisión conceptual}
    La muestra no reemplaza la definición de la población: delimita qué datos
    se observan para aprender sobre ella.
  \end{alertblock}
\end{frame}

\begin{frame}{Caso de lectura 1/4 --- interpretación}
  \footnotesize
  \begin{block}{Resolución conceptual}
    Primero se fija la población objetivo; después se elige la muestra; por
    último se cuantifica la incertidumbre de generalizar.
  \end{block}
  \begin{exampleblock}{Resultado}
    Una muestra grande puede reducir el error muestral, pero no corrige por sí
    sola un sesgo de selección o de medición.
  \end{exampleblock}
\end{frame}

\begin{frame}{Caso de lectura 2/4 --- parámetro y estadístico}
  \footnotesize
  La media poblacional $\mu$ es fija y desconocida. La media muestral $\bar X$
  cambia cuando se repite el muestreo.
  \begin{block}{Pregunta}
    ¿Qué cantidad puede calcularse directamente con los datos?
  \end{block}
  La respuesta de la sección es el estadístico $\bar X$, que sirve para estimar
  $\mu$.
\end{frame}

\begin{frame}{Caso de lectura 2/4 --- estimación}
  \footnotesize
  \[
    \widehat{\mu}=\bar X.
  \]
  \begin{itemize}\itemsep=0.12cm
    \item La estimación puntual entrega un valor.
    \item La estimación por intervalo añade un rango y un nivel de confianza.
    \item Ambos procedimientos dependen de la distribución muestral.
  \end{itemize}
\end{frame}

\begin{frame}{Caso de lectura 3/4 --- proceso de inferencia}
  \footnotesize
  \begin{enumerate}\itemsep=0.13cm
    \item Definir la característica.
    \item Seleccionar la muestra.
    \item Calcular el estadístico.
    \item Elegir el procedimiento inferencial.
    \item Interpretar en el contexto.
  \end{enumerate}
  Estos pasos son la traducción operativa de la lista de la nota de lectura.
\end{frame}

\begin{frame}{Caso de lectura 3/4 --- fuentes de error}
  \footnotesize
  \begin{block}{Resolución conceptual}
    El error de muestreo se modela mediante la distribución muestral; los
    errores no muestrales se reducen con diseño, medición y seguimiento.
  \end{block}
  \begin{alertblock}{Conclusión}
    Una inferencia válida requiere atender ambos tipos de error.
  \end{alertblock}
\end{frame}

\begin{frame}{Caso de lectura 4/4 --- papel del TLC}
  \footnotesize
  El TLC permite aproximar por una normal la distribución de una media muestral
  cuando se cumplen sus condiciones y $n$ es suficientemente grande.
  \begin{block}{Uso posterior}
    Esa aproximación sostiene pruebas Z/t e intervalos de confianza que se
    desarrollan en el capítulo.
  \end{block}
\end{frame}

\begin{frame}{Síntesis de la sección}
  \footnotesize
  \begin{itemize}\itemsep=0.12cm
    \item Población y muestra delimitan el problema.
    \item Parámetros y estadísticos distinguen objetivo y evidencia.
    \item Estimación y pruebas son las dos formas de inferencia.
    \item El error muestral y el TLC explican la incertidumbre.
  \end{itemize}
\end{frame}

\begin{frame}{Transición curricular}
  \footnotesize
  \begin{block}{Lo que queda establecido}
    Inferir significa pasar de estadísticos muestrales a afirmaciones sobre una
    población, registrando el error y las condiciones del procedimiento.
  \end{block}
  \begin{alertblock}{Siguiente sección}
    \textbf{05.01 Transformación de variables}: herramientas para obtener la
    distribución de $Y=g(X)$ y preparar el estudio de distribuciones muestrales.
  \end{alertblock}
\end{frame}
\end{document}
